\section{The Effects of the Y220C Mutation and PK11000 on p53 Isoforms}
\label{chap:appendix:isoforms}

\begin{figure}
\centering
\includegraphics[width=\textwidth]{Isoforms/p53_Isoforms_DBDp_RMSD_vs_Time.png}
\caption[RMSD vs Time of p53 Isoforms (DBD')]{Root Mean Square Deviation of the DBD' subset of each isoform over time. WT shown in blue, WT\_PK11000 shown in yellow, Y220C shown in red, and Y220C\_PK11000 shown in green. Averages over three replicates are shown in bold with a ghost representing the standard error of each frame.}
\label{fig:isoforms:rmsd-vs-time-DBDp}
\end{figure}

\begin{figure}
\centering
\includegraphics[width=\textwidth]{Isoforms/p53_Isoforms_DBDpp_RMSD_vs_Time.png}
\caption[RMSD vs Time of p53 Isoforms (DBD'')]{Root Mean Square Deviation of the DBD'' subset of each isoform over time. WT shown in blue, WT\_PK11000 shown in yellow, Y220C shown in red, and Y220C\_PK11000 shown in green. Averages over three replicates are shown in bold with a ghost representing the standard error of each frame.}
\label{fig:isoforms:rmsd-vs-time-DBDpp}
\end{figure}

\begin{figure}
\centering
\includegraphics[width=\textwidth]{Isoforms/p53_Isoforms_FL_RoG_vs_Time.png}
\caption[RoG vs Time of p53 Isoforms]{Radius of Gyration of each isoform over time. WT shown in blue, WT\_PK11000 shown in yellow, Y220C shown in red, and Y220C\_PK11000 shown in green. Averages over three replicates are shown in bold with a ghost representing the standard error of each frame.}
\label{fig:isoforms:rog-vs-time}
\end{figure}

\begin{figure}
\centering
\includegraphics[width=\textwidth]{Isoforms/p53_Isoforms_DBDp_RoG_vs_Time.png}
\caption[RoG vs Time of p53 Isoforms (DBD')]{Radius of Gyration of the DBD' of each isoform over time. WT shown in blue, WT\_PK11000 shown in yellow, Y220C shown in red, and Y220C\_PK11000 shown in green. Averages over three replicates are shown in bold with a ghost representing the standard error of each frame.}
\label{fig:isoforms:rog-vs-time-DBDp}
\end{figure}

\begin{figure}
\centering
\includegraphics[width=\textwidth]{Isoforms/p53_Isoforms_DBDpp_RoG_vs_Time.png}
\caption[RoG vs Time of p53 Isoforms (DBD'')]{Radius of Gyration of the DBD'' of each isoform over time. WT shown in blue, WT\_PK11000 shown in yellow, Y220C shown in red, and Y220C\_PK11000 shown in green. Averages over three replicates are shown in bold with a ghost representing the standard error of each frame.}
\label{fig:isoforms:rog-vs-time-DBDpp}
\end{figure}

\begin{figure}
\centering
\includegraphics[width=\textwidth]{Isoforms/p53_Isoforms_RoG_Violin_DBDp.png}
\caption[p53 Isoform RMSD Violin Plot for DBD' Subset]{A violin plot showing the distribution of RMSD values for the DBD' subset of each isoform. WT shown in blue, Y220C shown in red, and Y220C\_PK11000 shown in green. For each system, the median value of the distribution is shown with a horizontal line colored according to the system.}
\label{fig:isoforms:rmsd-violin-DBDp}
\end{figure}

\begin{figure}
\centering
\includegraphics[width=\textwidth]{Isoforms/p53_Isoforms_RoG_Violin_DBDpp.png}
\caption[p53 Isoform RMSD Violin Plot for DBD'' Subset]{A violin plot showing the distribution of RMSD values for the DBD'' subset of each isoform. WT shown in blue, Y220C shown in red, and Y220C\_PK11000 shown in green. For each system, the median value of the distribution is shown with a horizontal line colored according to the system.}
\label{fig:isoforms:rmsd-violin-DBDpp}
\end{figure}

\begin{figure}
\centering
\includegraphics[height=0.4\textheight]{Isoforms/p53_RoG_RescueScore.png}
\caption[p53 Isoform RoG Rescue Score]{The PK11000 rescue score, using the metric of Radius of Gyration applied to the full protein, the DBD', and DBD''. A rescue score between 0 and 1 indicates a partial rescue, a score greater than 1 indicates an overshoot of the WT, and a score less than 0 indicates that Y220C\_PK11000 is further from the WT than Y220C.}
\label{fig:isoforms:rog-rescue-score}
\end{figure}

\begin{figure}
\centering
\includegraphics[width=\textwidth]{Isoforms/p53_Isoforms_EE_diff.png}
\caption[p53 Isoform \(\Delta\)(Embedding Error) Score]{The difference between Y220C\_PK11000 and Y220C local embedding error score for each residue for each isoform, relative to the WT. A value less than zero indicates PK11000 restored WT-like interactions compared to the Y220C, and a value greater than zero indicates the opposite. Secondary structure is shown for each residue; blue indicates no secondary structure, red indicates a \(\beta\)-sheet for that residue, and yellow indicates an \(\alpha\)-helix.}
\label{fig:isoforms:ee-diff}
\end{figure}